\subsubsection{Circuito RL en paralelo}
En esta parte de la experiencia se distingue un error poco menor al 10\% en la medición de la corriente que entrega la fuente comparado con las simulaciones realizadas(Véase la tabla \textbf{REFERENCIAR LA TABLA CON LOS DATOS MEDIDOS Y CALCULADOS DE RL PARALELO PARTE 2}). Este error resulta relativamente pequeño. Aun así, puede ser asociado con el comportamiento no ohmico de la resistencia. La resistencia real de este experimento son lámparas incandescentes que al calentarse aumentan su resistividad. El calor disipado por las lámparas es proporcional a la potencia disipada. Si observa la tabla \textbf{REFERENCIAR LA TABLA DE LA PARTE 1} puede notar que la corriente que circula solamente por la lámpara es de 600mA mientras que en esta experiencia la corriente entregada por la fuente es un 17\% mayor. A su vez, la corriente que entrega la fuente es la suma de las corrientes sobre la resistencia en paralelo y el inductor(con su resistencia asociada). Por lo tanto, se puede inferir que la potencia que entrega en forma de calor la resistencia es menor a la de la primera parte (cuyos datos tomamos para calcular la resistencia simulada) y por lo tanto representa una resistencia menor. Como la resistencia es menor a la tomada para las simulaciones es natural que se note una corriente mayor en los experimentos.

De la tabla \textbf{REFERENCIAR TABLA DE CALCULOS DE CIRCUITO RL PARALELO} se identifica una predominancia resistiva evidenciada por un factor de potencia alto y por una potencia activa sustancialmente mayor a la potencia reactiva. Se puede apreciar una impedancia mayor en compararación con el circuito puramente resistivo pues se le agrega al circuito la impedancia de la bobina y de su resistencia interna. 


